\section{Answer Pools}
\label{app:data}

\begin{table}[h]
\centering
\small
\begin{tabular}{llrr}
\toprule
Benchmark & Answering model & Answers & Right \\
\midrule
MMLU-Redux & Qwen3.5-4B & 1,652 & 43\% \\
 & Qwen3.6-35B-A3B & 922 & 81\% \\
 & Nemotron-3-Super-120B & 922 & 72\% \\
\midrule
MATH & Qwen3.5-4B & 1,194 & 50\% \\
 & Qwen3.6-35B-A3B & 602 & 86\% \\
 & Nemotron-3-Super-120B & 588 & 75\% \\
\midrule
BBEH & 8 models & 2,244 & 34--58\% \\
\bottomrule
\end{tabular}
\caption{Main answer pools. BBEH answers come from Qwen3.5-4B, Qwen3.5-9B, Qwen3.6-35B-A3B, Nemotron-3-Nano-30B, Nemotron-3-Super-120B, DeepSeek-V3.1, gpt-oss-20b and gpt-oss-120b. All answers end with a final answer line; MMLU-Redux answers are graded by their final letter against the expert key, MATH answers by GLM-5.3-Flash against the reference (exact matches skip the grader), BBEH answers by the benchmark's matching rule.}
\label{tab:data}
\end{table}

\section{Judge Prompts}
\label{app:prompts}

All judges see one user turn. For MMLU-Redux (MATH and BBEH differ only in the first line: ``a math problem'', ``a reasoning question'', with \texttt{[PROBLEM]} for MATH):

\begin{quote}\small\ttfamily
You are grading an AI assistant's answer to a multiple-choice question.\\[2pt]
[QUESTION]\\ \{question\}\\[2pt]
[ASSISTANT'S ANSWER]\\ \{answer\}\\[2pt]
Is the assistant's final answer correct? Only the final answer matters, not the style of the reasoning. Reason briefly, then end with exactly one line: "VERDICT: CORRECT" or "VERDICT: INCORRECT".
\end{quote}

For the forced verdict and all hidden-state measurements, the assistant turn starts with \texttt{VERDICT:} and we read the next token. For the reasoning baseline, the judge generates freely (greedy, stopping at \texttt{VERDICT:}; at most 8,192 tokens in a 32K window for BBEH and for Llama-3.1-8B on MATH, 1,536 tokens in a 16K window otherwise) and we read the next token after appending \texttt{VERDICT:} to its reasoning. Runs that never write a verdict within the budget (mostly Qwen3.5-4B and Llama-3.1-8B on BBEH, 2--8\% of their answers) are repetitive loops; they are scored the same way. The verdict score uses the first tokens of `` CORRECT'' and `` INCORRECT''.

\section{Layer Rule and Fresh Test}
\label{app:rule}

\paragraph{Rule, fixed before the main run.} Gap check: stop if the gap is $\le 5$. Push test: on 400 random answers, the accept rate at $-8$ and $+8$ standard deviations along the direction at every second layer; a layer responds if the accept rate rises by at least 10 points. Layer: the best-ranking responding layer; if none responds, the best-ranking layer. Stretch: $\times$11 at that layer only, along the direction and along a random direction; $\times$4 and $\times$21 are recorded as secondary readouts.

\paragraph{Revision after a failure.} Mistral-7B on MMLU-Redux accepts 99\% of answers, so no push changes its decision, although the verdict score follows the direction (within-question correlation $+0.54$ to $+0.85$ from layer 18). We therefore measure the response on the verdict score: sensitivity is the mean change in verdict score per standard deviation pushed, divided by the verdict score's within-question standard deviation, and a layer responds if its sensitivity is at least 0.05. In a full sweep this picks layer 16 for Mistral-7B on MMLU-Redux ($+2.3$ [$+0.5$, $+3.7$]). The seven original judges were stretched at the layer the original rule chose (for Mistral-7B on MMLU-Redux, layer 4: $-0.7$), and all pooled numbers use those layers; the revised rule chose the layers for Llama-3.2-1B, Qwen2.5-7B and Llama-3.1-70B. In the main text and Table~\ref{tab:pairs}, a pair counts as responding if the verdict's sensitivity at the stretched layer is at least 0.05.

\paragraph{Fresh test.} Before running Llama-3.2-1B and Qwen2.5-7B, new to the study, we wrote down the revised rule and the prediction that every pair with a gap and a responding layer gains with $\times$11 (interval above 0) and that the random direction stays within $\pm1$. It held for Llama-3.2-1B on MMLU-Redux ($+5.4$) and MATH ($+7.9$), and the random direction stayed within $\pm1$ throughout. It failed on BBEH for both judges ($-1.5$ and $+1.0$). In both, the direction's ranking is flat across the responding layers to within split noise, and the rule took the layer where the verdict only begins to respond; a neighbouring layer with a stronger response gains $+6.4$ [$+3.3$, $+10.3$] (Llama-3.2-1B, layer 14) and $+4.6$ [$+1.2$, $+8.2$] (Qwen2.5-7B, layer 22). Qwen2.5-7B on MMLU-Redux has no responding layer; stretching its best-ranking layer still gains $+2.5$ [$+1.1$, $+4.0$].

\paragraph{Strength.} Judges whose verdict responds strongly (sensitivity $\ge 0.2$) do best at $\times$4--$\times$11; weakly responding ones (sensitivity $< 0.08$, mostly BBEH) at $\times$11--$\times$21. $\times$11 is within about two points of the best tested strength for most pairs.

\section{Every Judge and Benchmark}
\label{app:pairs}

\begin{table*}[h]
\centering
\small
\begin{tabular}{ll rr rr rrrr}
\toprule
Bench & Judge & Gap & Gap, 20 splits & Layer & Sens. & Verdict & Reason & Stretch & Direction \\
\midrule
MMLU-Redux & Llama-3.1-8B & +12.6 & +12.2 $\pm$ 0.8 & 16 & 0.13 & 51.1 & 52.6 & 58.0 & 63.7 \\
MMLU-Redux & Llama-3.2-3B & +11.6 & -- & 12 & 0.06 & 51.7 & 46.5 & 55.8 & 63.4 \\
MMLU-Redux & Qwen3.5-9B & +0.8 & -- & \multicolumn{6}{l}{gap $\le$ 5: not stretched} \\
MMLU-Redux & Qwen3.5-4B & +1.2 & -- & \multicolumn{6}{l}{gap $\le$ 5: not stretched} \\
MMLU-Redux & Phi-4-mini & +7.2 & +7.8 $\pm$ 0.4 & 16 & 0.12 & 55.0 & 52.7 & 58.9 & 61.1 \\
MMLU-Redux & Mistral-7B & +7.6 & +7.9 $\pm$ 0.3 & 4$^*$ & -0.00 & 51.5 & 51.3 & 50.8 & 59.1 \\
MMLU-Redux & Granite-3.3-8B & +6.4 & +5.9 $\pm$ 0.5 & 20 & 0.11 & 55.4 & 51.6 & 57.2 & 61.9 \\
MMLU-Redux & Llama-3.2-1B & +12.1 & +11.7 $\pm$ 0.5 & 10 & 0.05 & 49.6 & 50.9 & 55.0 & 59.6 \\
MMLU-Redux & Qwen2.5-7B & +10.3 & +10.0 $\pm$ 0.7 & 20$^*$ & 0.04 & 51.8 & 50.0 & 54.3 & 62.2 \\
MMLU-Redux & Llama-3.1-70B & +1.7 & +1.4 $\pm$ 0.2 & 34 & 0.27 & 65.0 & 62.6 & 64.9 & 66.8 \\
\midrule
MATH & Llama-3.1-8B & +13.0 & -- & 16 & 0.26 & 54.3 & 53.0 & 62.5 & 66.4 \\
MATH & Llama-3.2-3B & +11.8 & +11.5 $\pm$ 0.4 & 14 & 0.26 & 53.5 & 52.3 & 59.8 & 65.1 \\
MATH & Qwen3.5-9B & +4.9 & -- & \multicolumn{6}{l}{gap $\le$ 5: not stretched} \\
MATH & Qwen3.5-4B & +2.0 & +2.2 $\pm$ 0.3 & \multicolumn{6}{l}{gap $\le$ 5: not stretched} \\
MATH & Phi-4-mini & +12.9 & +12.5 $\pm$ 0.4 & 14 & 0.12 & 55.9 & 56.2 & 62.1 & 68.7 \\
MATH & Mistral-7B & +5.7 & +5.9 $\pm$ 0.7 & 16 & 0.20 & 53.4 & 53.1 & 55.5 & 57.8 \\
MATH & Granite-3.3-8B & +4.4 & +4.5 $\pm$ 0.3 & \multicolumn{6}{l}{gap $\le$ 5: not stretched} \\
MATH & Llama-3.2-1B & +12.8 & +12.1 $\pm$ 1.1 & 14 & 0.15 & 51.3 & 52.5 & 59.1 & 64.3 \\
MATH & Qwen2.5-7B & +5.0 & +5.0 $\pm$ 0.3 & \multicolumn{6}{l}{gap $\le$ 5: not stretched} \\
MATH & Llama-3.1-70B & +4.4 & +4.5 $\pm$ 0.3 & 34 & 0.25 & 61.8 & 62.5 & 62.3 & 66.3 \\
\midrule
BBEH & Llama-3.1-8B & +11.2 & +9.2 $\pm$ 2.3 & 14$^*$ & 0.04 & 55.1 & 50.8 & 61.8 & 66.3 \\
BBEH & Llama-3.2-3B & +10.3 & +8.8 $\pm$ 3.1 & 14$^*$ & 0.03 & 50.9 & 51.4 & 55.6 & 60.9 \\
BBEH & Qwen3.5-9B & -1.5 & -- & \multicolumn{6}{l}{gap $\le$ 5: not stretched} \\
BBEH & Qwen3.5-4B & +9.3 & +8.0 $\pm$ 1.6 & 20 & 0.08 & 59.5 & 63.9 & 64.7 & 68.7 \\
BBEH & Phi-4-mini & +10.8 & +7.2 $\pm$ 4.2 & 20 & 0.07 & 52.0 & 46.9 & 56.0 & 62.7 \\
BBEH & Mistral-7B & +4.5 & +5.0 $\pm$ 0.8 & \multicolumn{6}{l}{gap $\le$ 5: not stretched} \\
BBEH & Granite-3.3-8B & +7.4 & +6.5 $\pm$ 1.6 & 22 & 0.19 & 58.8 & 53.1 & 65.1 & 66.2 \\
BBEH & Llama-3.2-1B & +7.9 & +5.9 $\pm$ 1.5 & 10 & 0.20 & 49.2 & 52.4 & 47.7 & 57.1 \\
BBEH & Qwen2.5-7B & +10.1 & +8.5 $\pm$ 1.2 & 20 & 0.09 & 53.2 & 54.6 & 54.2 & 63.4 \\
BBEH & Llama-3.1-70B & +6.0 & +4.6 $\pm$ 2.4 & 32 & 0.06 & 58.6 & 57.9 & 60.0 & 64.6 \\
\bottomrule
\end{tabular}

\caption{Every judge--benchmark pair. Gap: best layer's direction ranking minus the verdict's, on the split used for the layer rule; 20 splits: mean $\pm$ standard deviation over 20 random half-splits of the questions (at the best layer). Layer: the stretched layer ($^*$: the verdict's sensitivity there is below 0.05). Sens.: the verdict's sensitivity there. Llama-3.1-70B is the scale check of \S\ref{sec:scale}. Last four columns: ranking ($\times$100) of the forced verdict, the verdict after reasoning, the verdict after stretching $\times$11, and the direction, on the main answer pool.}
\label{tab:pairs}
\end{table*}

\begin{table*}[h]
\centering
\small
\begin{tabular}{lll r rrrrrr}
\toprule
Bench & Judge & Answers by & Groups & Verdict & Reason & Str.\,$\times$4 & Str.\,$\times$11 & Direction & Verd.+dir. \\
\midrule
MATH & Llama-3.1-70B & OLMo-2-7B & 89 & 62.6 & 62.8 & 62.6 & 58.9 & 64.4 & 63.3 \\
MATH & Llama-3.1-70B & Qwen3.5-4B & 280 & 62.9 & 63.3 & 63.7 & 63.1 & 64.0 & 64.6 \\
MATH & Llama-3.1-8B & OLMo-2-7B & 154 & 58.9 & 60.5 & 58.6 & 57.9 & 58.9 & 59.0 \\
MATH & Llama-3.1-8B & Qwen3.5-4B & 350 & 58.3 & 57.4 & 62.0 & 62.2 & 60.5 & 61.8 \\
MATH & Llama-3.1-8B & Qwen3.5-4B, failed 9/9 & 255 & 56.8 & 56.5 & 61.6 & 61.8 & 59.9 & 61.6 \\
MATH & Llama-3.2-1B & OLMo-2-7B & 150 & 54.8 & 54.6 & 54.4 & 54.3 & 53.3 & 52.5 \\
MATH & Llama-3.2-1B & Qwen3.5-4B & 340 & 56.2 & 54.4 & 58.3 & 59.9 & 61.8 & 59.9 \\
MATH & Llama-3.2-3B & OLMo-2-7B & 132 & 60.5 & 58.4 & 60.2 & 60.0 & 56.8 & 60.1 \\
MATH & Llama-3.2-3B & Qwen3.5-4B & 327 & 62.7 & 58.8 & 65.3 & 59.4 & 61.3 & 65.3 \\
MATH & Mistral-7B & OLMo-2-7B & 142 & 58.2 & -- & 58.6 & 56.6 & 56.3 & 57.9 \\
MATH & Mistral-7B & Qwen3.5-4B & 340 & 55.8 & -- & 56.4 & 55.2 & 58.2 & 57.8 \\
MATH & Phi-4-mini & OLMo-2-7B & 102 & 67.0 & 61.4 & 66.0 & 64.7 & 53.3 & 62.7 \\
MATH & Phi-4-mini & Qwen3.5-4B & 307 & 58.5 & 57.7 & 60.2 & 56.4 & 61.5 & 59.5 \\
BBEH & Granite-3.3-8B & 4 pool models & 335 & 50.5 & 51.3 & 53.5 & 53.9 & 56.0 & 54.6 \\
BBEH & Llama-3.1-70B & 4 pool models & 311 & 55.2 & 58.0 & 56.1 & 54.7 & 56.1 & 57.8 \\
BBEH & Llama-3.1-8B & 4 pool models & 303 & 54.3 & 50.1 & 54.6 & 55.4 & 55.7 & 56.6 \\
BBEH & Llama-3.2-1B & 4 pool models & 350 & 53.7 & 53.5 & 54.2 & 49.8 & 55.3 & 54.3 \\
BBEH & Llama-3.2-3B & 4 pool models & 341 & 53.2 & 54.1 & 55.5 & 54.5 & 55.6 & 54.9 \\
BBEH & Phi-4-mini & 4 pool models & 343 & 55.7 & 53.8 & 55.9 & 56.6 & 58.7 & 57.3 \\
BBEH & Qwen3.5-4B & 4 pool models & 210 & 54.4 & 56.2 & 54.8 & 55.2 & 57.9 & 56.2 \\
MMLU-Redux & Granite-3.3-8B & OLMo-2-7B & 193 & 47.7 & 50.4 & 48.8 & 48.7 & 50.6 & 49.1 \\
MMLU-Redux & Granite-3.3-8B & Qwen3.5-4B & 271 & 44.9 & 48.6 & 44.3 & 45.1 & 46.3 & 44.9 \\
MMLU-Redux & Llama-3.1-70B & OLMo-2-7B & 129 & 50.6 & 46.2 & 49.4 & 48.7 & 50.5 & 49.9 \\
MMLU-Redux & Llama-3.1-70B & Qwen3.5-4B & 183 & 49.4 & 50.4 & 49.5 & 46.0 & 48.9 & 49.8 \\
MMLU-Redux & Llama-3.1-8B & OLMo-2-7B & 305 & 53.3 & 52.7 & 53.5 & 53.9 & 53.3 & 53.6 \\
MMLU-Redux & Llama-3.1-8B & Qwen3.5-4B & 401 & 47.5 & 50.3 & 47.8 & 47.5 & 51.0 & 47.4 \\
MMLU-Redux & Llama-3.2-1B & OLMo-2-7B & 241 & 51.5 & 54.6 & 51.1 & 49.5 & 49.3 & 51.9 \\
MMLU-Redux & Llama-3.2-1B & Qwen3.5-4B & 318 & 50.0 & 49.8 & 49.9 & 48.1 & 50.5 & 50.0 \\
MMLU-Redux & Phi-4-mini & OLMo-2-7B & 198 & 50.5 & 53.0 & 50.5 & 50.8 & 51.9 & 52.2 \\
MMLU-Redux & Phi-4-mini & Qwen3.5-4B & 269 & 43.6 & 45.0 & 45.2 & 45.3 & 45.8 & 45.2 \\
\bottomrule
\end{tabular}

\caption{Same-model pairs on questions the judge cannot solve, per judge: ranking ($\times$100) by each readout. Groups: questions with at least one right and one wrong answer from the same answering model.}
\label{tab:smpairs}
\end{table*}

\section{Authorship Directions and Controls}
\label{app:auth}

Per judge and benchmark, at the best-ranking layer and with the same two question halves as the correctness direction, we take, within every question and correctness class that holds answers by both models, the difference between the two models' mean states, and average these paired differences: Qwen3.6-35B-A3B vs Qwen3.5-4B (``own pair'') and Nemotron-3-Super-120B vs Qwen3.5-4B. Table~\ref{tab:auth} reports, on held-out questions, how well the first direction separates answers by its own pair and by the other pair (within question and correctness class; 50 = none, below 50 = reversed), its cosine with the correctness direction, and the ranking of right above wrong answers by the verdict, by the correctness direction, and by the correctness direction with both authorship directions projected out ($\perp$). Llama-3.1-8B on MATH and Llama-3.2-3B on MMLU-Redux were run with a different storage layout and are not in this table; the controls below include them. The Verdict column is recomputed on all answers with stored hidden states and can differ by a few tenths from Table~\ref{tab:pairs}, which uses only questions where every readout has a score.

\paragraph{Controls.} All fits use the training half of the questions; numbers in \S\ref{sec:what} are pooled over the \R{ctl.both.n} main-pool pairs on MMLU-Redux and MATH (\R{ctl.mmlu.n} and \R{ctl.math.n}). \emph{Label shuffles}: within every question and correctness class of the training half, the answering-model labels are permuted, the two paired directions are fitted as before, and their span is projected out of the correctness direction; 100 shuffles per judge. The mean removal over shuffles is at most \R{ctl.mmlu.perm.maxmean} points per judge on MMLU-Redux and \R{ctl.math.perm.maxmean} on MATH (it is often negative: removing two high-variance directions can sharpen the ranking). \emph{Random subspaces}: 100 isotropic random 2-D subspaces per judge, drawn without labels, remove at most \R{ctl.mmlu.iso.max} points. \emph{Within-model fit}: the correctness direction is fitted after centring each question's answers separately for each answering model, so only right-vs-wrong contrasts between answers by the same model enter. \emph{Held-out answering model}: the direction is fitted without one model's answers and scored on that model's same-model pairs; only Qwen3.5-4B has several answers per question, so only it can be held out this way (fitted without it: \R{ctl.ho.c_ho.gain.ci} over the verdict; with it: \R{ctl.ho.c_all.gain.ci}). OLMo-2 in \S\ref{sec:beyond} is a fully held-out answering model.

\paragraph{Why not a variance share.} A simpler measure of authorship is the share of a score's within-question variance that lies between answering models (\R{share.direct.lo}--\R{share.direct.hi}\% for the direction, \R{share.verdict.lo}--\R{share.verdict.hi}\% for the verdict). With one or two answers per model and question, a score shuffled within question already gets \R{share.null.lo}--\R{share.null.hi}\%, and the correctness label itself \R{share.gold.lo}--\R{share.gold.hi}\%, so this share says little about authorship. The pair split and the authorship directions of \S\ref{sec:what} replace it.

\begin{table*}[h]
\centering
\small
\begin{tabular}{ll r rr r rrr rr}
\toprule
& & & \multicolumn{2}{c}{Authorship dir.} & & \multicolumn{3}{c}{Ranking, all pairs} & \multicolumn{2}{c}{Same-model pairs} \\
\cmidrule(lr){4-5}\cmidrule(lr){7-9}\cmidrule(l){10-11}
Bench & Judge & Layer & Own pair & Other pair & cos & Verdict & Dir. & Dir.$\perp$ & Dir. & Dir.$\perp$ \\
\midrule
MMLU-Redux & Llama-3.1-8B & 16 & 78.7 & 30.9 & +0.19 & 51.1 & 63.7 & 58.4 & 60.7 & 59.3 \\
MMLU-Redux & Phi-4-mini & 14 & 79.9 & 32.6 & +0.13 & 55.0 & 62.2 & 56.9 & 54.3 & 55.6 \\
MMLU-Redux & Mistral-7B & 4 & 80.8 & 27.4 & $-$0.05 & 51.5 & 59.1 & 52.7 & 51.5 & 49.4 \\
MMLU-Redux & Granite-3.3-8B & 20 & 85.1 & 29.8 & +0.06 & 55.4 & 61.8 & 58.4 & 57.6 & 57.9 \\
MMLU-Redux & Llama-3.2-1B & 4 & 84.3 & 28.5 & +0.14 & 49.6 & 61.8 & 53.5 & 52.8 & 52.8 \\
MMLU-Redux & Qwen2.5-7B & 20 & 77.6 & 31.4 & +0.19 & 51.8 & 62.1 & 57.6 & 57.0 & 57.6 \\
\midrule
MATH & Llama-3.2-3B & 14 & 84.4 & 42.3 & +0.43 & 53.7 & 65.5 & 62.8 & 61.6 & 62.1 \\
MATH & Qwen3.5-4B & 20 & 66.2 & 52.7 & +0.71 & 66.4 & 68.4 & 69.5 & 64.4 & 68.2 \\
MATH & Phi-4-mini & 14 & 71.3 & 45.5 & +0.76 & 56.0 & 68.9 & 67.1 & 65.2 & 63.1 \\
MATH & Mistral-7B & 2 & 73.1 & 91.2 & +0.52 & 53.4 & 59.1 & 60.8 & 48.2 & 59.2 \\
MATH & Granite-3.3-8B & 22 & 77.1 & 44.3 & +0.68 & 62.5 & 66.9 & 62.3 & 59.1 & 58.5 \\
MATH & Llama-3.2-1B & 14 & 81.0 & 35.3 & +0.44 & 51.9 & 64.7 & 62.3 & 58.9 & 58.7 \\
MATH & Qwen2.5-7B & 20 & 69.0 & 43.5 & +0.57 & 58.5 & 63.5 & 64.4 & 60.0 & 60.7 \\
\bottomrule
\end{tabular}

\caption{Authorship directions (held out). Own pair / other pair: AUROC of the Qwen3.6-vs-Qwen3.5-4B direction for Qwen3.6 vs Qwen3.5-4B answers and for Nemotron vs Qwen3.5-4B answers of equal correctness. cos: with the correctness direction. Ranking of right above wrong answers by the verdict, the correctness direction, and the correctness direction with both authorship directions removed ($\perp$), on all and on same-model pairs.}
\label{tab:auth}
\end{table*}

\section{Best-of-$n$ Selection}
\label{app:bon}

\begin{table*}[h]
\centering
\small
\begin{tabular}{lr rr rrrr}
\toprule
Answers & Judges & Random & Verdict & Reason first & Stretch $\times$11 & Direction & Verdict + dir. \\
\midrule
All questions & 19 & 59.7 & 61.6 & $-$1.7 & \textbf{+2.5} & \textbf{+8.1} & \textbf{+5.4} \\
Judge cannot solve & 19 & 58.4 & 58.3 & $-$1.0 & \textbf{+2.4} & \textbf{+9.0} & \textbf{+5.7} \\
Same model, MATH, Qwen3.5-4B & 5 & 56.3 & 63.6 & $-$4.7 & +0.9 & +0.6 & +1.1 \\
Same model, MATH, OLMo-2-7B & 5 & 36.0 & 47.1 & $-$1.5 & $-$1.5 & $-$3.6 & $-$2.8 \\
Same model, BBEH, 4 pool models & 6 & 47.5 & 50.3 & $-$0.9 & \textbf{+1.5} & \textbf{+2.3} & \textbf{+1.9} \\
Same model, MMLU-Redux, Qwen3.5-4B & 4 & 48.2 & 44.4 & +2.1 & $-$1.2 & +0.2 & $-$1.7 \\
Same model, MMLU-Redux, OLMo-2-7B & 4 & 43.2 & 41.7 & +2.3 & $-$0.5 & +1.4 & +0.7 \\
\bottomrule
\end{tabular}

\caption{Best-of-$n$ accuracy: the share of groups (questions; question $\times$ answering model for same-model sets; only groups with at least one right and one wrong answer) in which the top-scored answer is right, for random choice and the forced verdict, and the gain of each readout over the verdict (mean over judges; bold: 95\% interval above 0).}
\label{tab:bon}
\end{table*}

\section{Pooling}
\label{app:pool}

The primary estimate is the equal-weight mean over judge--benchmark pairs, with questions resampled within benchmark (2,000 draws) and the same draws for every judge and readout. An earlier version weighted pairs by the inverse variance of their own intervals and treated judges as independent; that weighting favours pairs with tight intervals, which are often those with small gains, and its intervals are too narrow because judges grade the same answers. On the main pools it gives \R{legacy.pool.direct.rank.ci} for the direction, \R{legacy.pool.sum.rank.ci} for verdict + direction, \R{legacy.pool.x11.rank.ci} for stretching and \R{legacy.pool.cot.rank.ci} for reasoning first, against \R{pool.direct.rank.ci}, \R{pool.sum.rank.ci}, \R{pool.x11.rank.ci} and \R{pool.cot.rank.ci} here. Weighting every question equally (the mean of per-question AUROCs) instead of every right--wrong pair gives \R{pool.direct.rank_macro.ci}, \R{pool.sum.rank_macro.ci}, \R{pool.x11.rank_macro.ci} and \R{pool.cot.rank_macro.ci}. All of these are conditional on the judges and benchmarks studied; they do not generalise to a population of judges.

\section{Surface Features}
\label{app:surface}

Features of the answer text (the block the judge grades): log characters, log lines, count of bold markers, count of \texttt{\$}, count of the word ``step'', counts of hedge words (maybe, might, perhaps, possibly, probably, likely, approximately, I think, I believe, not sure, unclear, assume), confident words (clearly, definitely, certainly, obviously, indeed, undoubtedly, must be) and self-correction words (wait, actually, however, recheck, double-check, let me re-, mistake, correction), and whether the text contains ``the answer is''. Features, labels and scores are centred within question and the features standardised; a least-squares fit on one half of the questions predicts the other half. \emph{Surface}: correctness predicted from the features. \emph{Direction without surface}: the direction score minus its prediction from the features. On same-model pairs the surface model ranks at \R{surf.surface.level.same} (verdict \R{surf.verdict.level.same}, direction \R{surf.direct.level.same}). On MMLU-Redux the cue is length rather than bare lines: \R{fmt.mmlu.nemo.0.bare}\% of Nemotron's wrong and \R{fmt.mmlu.nemo.1.bare}\% of its right answers are bare final lines, but \R{fmt.mmlu.q36.0.long}\% of Qwen3.6's wrong and \R{fmt.mmlu.q36.1.long}\% of its right answers exceed 8,000 characters. In the format-trimmed pools (answers of 100--8,000 characters) the direction gains \R{surf.trim.direct.gain.same} on same-model and \R{surf.trim.direct.gain.diff} on different-model pairs.

\section{Held-Out Evaluation of the Recipe}
\label{app:outer}

The scope was fixed before running: Llama-3.1-8B, Qwen2.5-7B, Phi-4-mini and Granite-3.3-8B on MMLU-Redux, MATH and BBEH, \R{outer.npairs} pairs, including two that the main pool leaves out because their in-sample gap is at most 5 points (\R{outer.notpaper}). The outer folds are the paper's two question halves; each is the test set once. On the training half only: (1)~the direction at every candidate layer (class-mean difference after centring by question); (2)~an inner split of the training questions gives each layer's cross-fitted ranking and the verdict's; the best layer is the readout layer, and the pair is eligible if it beats the verdict by more than 5 points; (3)~the push test ($\pm$8 standard deviations on 400 training answers, no labels) gives each layer's sensitivity, and the stretch layer is the best-ranking layer with sensitivity $\ge 0.05$, or the best-ranking layer if none responds. On the test half: the direction at the readout layer, the verdict stretched ($\times$11) along the training direction at every candidate layer, and a random direction stretched at the chosen layer. The rule takes the stretched verdict if the fold is eligible and the verdict otherwise. \R{outer.folds.eligible} of the \R{outer.folds} folds were eligible (Table~\ref{tab:outer}). Picking the stretch layer after seeing the test results would add \R{outer.optimism.mean} point on average.

\begin{table*}[t]
\centering
\small
\begin{tabular}{ll cc r rrr rr}
\toprule
& & & & & \multicolumn{3}{c}{Held out: gain over verdict} & \multicolumn{2}{c}{Paper (in-sample choice)} \\
\cmidrule(lr){6-8}\cmidrule(l){9-10}
& Judge & Eligible & Layer & Verdict & Direction & Rule & Random & Direction & Stretch \\
\midrule
MMLU-Redux & Llama-3.1-8B & y/y & 16/16 & 51.1 & +12.1 & +7.0 & +0.2 & +12.6 & +6.9 \\
MMLU-Redux & Phi-4-mini & y/y & 14/18 & 55.0 & +7.2 & +3.8 & $-$0.1 & +6.1 & +3.9 \\
MMLU-Redux & Granite-3.3-8B & y/y & 20/20 & 55.4 & +6.5 & +1.7 & +0.1 & +6.5 & +1.7 \\
MMLU-Redux & Qwen2.5-7B & y/y & 20/20 & 51.8 & +10.3 & +2.5 & 0.0 & +10.4 & +2.5 \\
\midrule
MATH & Llama-3.1-8B & y/y & 20/16 & 54.4 & +12.3 & +8.8 & $-$0.4 & +12.2 & +8.2 \\
MATH & Phi-4-mini & y/y & 14/16 & 56.0 & +10.2 & +5.1 & +0.1 & +12.8 & +6.2 \\
MATH & Granite-3.3-8B & n/n & 20/20 & 62.5 & +4.0 & 0.0 & +0.1 & -- & -- \\
MATH & Qwen2.5-7B & y/y & 20/20 & 58.5 & +8.0 & $-$1.4 & $-$0.3 & -- & -- \\
\midrule
BBEH & Llama-3.1-8B & n/y & 18/16 & 55.1 & +7.7 & +3.5 & $-$0.2 & +11.2 & +6.8 \\
BBEH & Phi-4-mini & y/y & 10/20 & 52.0 & +6.4 & +2.7 & +0.4 & +10.8 & +4.0 \\
BBEH & Granite-3.3-8B & n/y & 38/24 & 58.8 & +3.2 & +1.3 & $-$0.8 & +7.4 & +6.3 \\
BBEH & Qwen2.5-7B & n/y & 22/22 & 53.2 & +8.6 & +2.6 & +0.7 & +10.1 & +1.0 \\
\bottomrule
\end{tabular}

\caption{Held-out evaluation of the recipe. Eligible and Layer: whether the inner gap exceeded 5 points, and the chosen stretch layer, in each of the two folds. Verdict: ranking ($\times$100). Held out: gains over the verdict of the direction, the rule (stretch if eligible, else the verdict) and a random direction stretched at the chosen layer. Paper: gains with the layer chosen on the evaluated questions (--: not stretched in the main pool).}
\label{tab:outer}
\end{table*}

\section{Compute}

Hidden-state capture, push tests and stretching ran on single H100 GPUs with Hugging Face Transformers (one forward pass per answer, plus a few edited re-reads of the last token, which can reuse the key-value cache); about one GPU-hour per judge--benchmark pair. Llama-3.1-70B ran on two H200 GPUs on a university cluster (2--3 hours per benchmark for capture, push test and stretching; about 4 hours of vLLM for its own answers and the reasoning baseline). Reasoning judgments and OLMo-2 answers used vLLM. Directions and readouts are computed on CPU from stored hidden states.
